\section{Caveats}\label{sec:caveats}

\paragraph{The limit depends on the window model.}
\begin{itemize}
\item The window restores only disconnected vacuum fluctuations.
\item In $d=1$ the corrected descendant spacings still miss by amounts that do not decrease with $N$
 (Table~\ref{tab:d1-ident}).
\item In $d=2$ the $\eps$ expansion of the window model is correct only through third order: its fourth-order
 coefficient grows like $\ln N$ (Section~\ref{d2:sec}). In $d=3$ the fourth-order coefficient at $N=20$ is 60\% of
 $C_4(3)$ and still growing slowly; whether its $N\to\infty$ limit is wrong is not shown (Section~\ref{d3:sec}).
\end{itemize}

\paragraph{The method fails toward the short-range end in every dimension.}
\begin{itemize}
\item The tuned coupling is largest there.
\item In $d=1$ it also grows with the cutoff, and toward the short-range end $\Ds\to0$ the $\sigma$ gap that the
 tuning imposes, $E_\sigma-E_v=\Ds$, closes.
\item The failure sets in at $\Ds\le0.05$ in $d=1$, at $\Ds\le0.175$ in $d=2$, and near the endpoint in $d=3$.
\end{itemize}

\paragraph{Errors.}
\begin{itemize}
\item The quoted errors are half-spreads of the extrapolations in the cutoff, not standard deviations, and they do
 not include the dependence on the window prescription (Section~\ref{sec:d1-model}).
\item The exponents of Eq.~\eqref{eq:fitN} are fitted, not derived.
\item At $\Ds\ge0.15$ in $d=1$ the spreads are probably optimistic.
\item The classification \emph{final} versus \emph{one-sided} is a judgement based on the local exponent.
\end{itemize}

\paragraph{Order of the analysis.}
\begin{itemize}
\item The window was designed after its effect had been seen in the descendant spacings.
\item The fit protocol of Section~\ref{sec:d1-extrap} was carried over unchanged from the earlier \code{prod1--7}
 analysis.
\item The held-out test and the local exponent were added after the $N=40$ data had been seen.
\item The tail-model and prescription comparisons of Section~\ref{sec:d1-model} were made after the gap to NRG had
 been seen. For that reason they are given as shifts and are not turned into an error.
\end{itemize}

\paragraph{One dimension.}
\begin{itemize}
\item \emph{Strong coupling.} The $d=1$ values at $\Ds\le0.05$ are not limits. At 0.025 and below, the correction
 overshoots: the sequence turns over and falls without bound in $\ln N$.
\item \emph{Comparison with NRG.} Where the descriptions overlap ($0.075\le\Ds\le0.2$), HT lies above NRG at all six points, by up to
 $\HTq{C1}$. At $\Ds=0.1$ and 0.125 the gap exceeds the spread of our extrapolations. However, the uncorrected
 prescription extrapolates to limits $\HTq{M10}$ lower, below NRG, so the dependence on the prescription is as large
 as the gap (Section~\ref{sec:d1-model}).
\item \emph{Prescription dependence.} The quoted errors of Table~\ref{tab:d1} do not include the prescription
 dependence. The raw and window limits differ by $\HTq{M6}$ at $\Ds=0.125$, 0.1 and 0.075.
\end{itemize}

\input{caveats_d2}

\input{caveats_d3}
